% Generated draft: requires mathematical and source review.
\documentclass[11pt]{article}
\usepackage[margin=25mm]{geometry}
\usepackage{fontspec}
\setmainfont{Latin Modern Roman}
\usepackage{amsmath,amssymb,mathtools,unicode-math}
\setmathfont{Latin Modern Math}
\usepackage{enumitem,needspace,xurl,hyperref}
\setlength{\parindent}{0pt}
\setlength{\parskip}{5pt}
\setlength{\emergencystretch}{3em}
\newcommand{\N}{\mathbb{N}}
\newcommand{\Z}{\mathbb{Z}}
\newcommand{\Q}{\mathbb{Q}}
\newcommand{\R}{\mathbb{R}}
\newcommand{\C}{\mathbb{C}}
\newcommand{\F}{\mathbb{F}}
\newcommand{\A}{\mathbb{A}}
\newcommand{\PP}{\mathbb{P}}
\newcommand{\E}{\mathbb{E}}
\newcommand{\eps}{\varepsilon}
\newcommand{\dd}{\,\mathrm{d}}
\begin{document}
\section*{UnsolvedMath ID: 2307050}
\textbf{Source problem code:} catalogue entry
\subsection{Definitions and mathematical statement}
For a connected Euclidean-open set $U\subset\R^n$, define the rectifiably accessible boundary set by
\[
 A_{\mathrm{rect}}(U)=\{\xi\in\partial U:\exists\gamma:[0,1]\to\R^n\text{ continuous},\
 \gamma([0,1))\subset U,\ \gamma(1)=\xi,\ \operatorname{length}(\gamma)<\infty\}.
\]
The source attributes to Brelot--Choquet the assertion that harmonic measure is carried by this set. In measure-theoretic terms, for each applicable interior pole $x$, this means $\omega_U^x(\partial U\setminus A_{\mathrm{rect}}(U))=0$, rather than a claim about the topological support being a subset of the accessible set.

The fine topology is the potential-theoretic topology generated to make superharmonic functions continuous; a finely open set is open in this topology. T. Lyons and B. Øksendal ask whether the analogous assertion holds for finely open sets and Keldysh measure: is the relevant boundary measure carried by boundary points accessible from inside by finite-length rectifiable paths?

The source does not define Keldysh measure, its pole and existence assumptions, which boundary is used for a finely open set, or whether the path/end-point topology remains Euclidean or is fine. It also does not say whether ordinary connectedness is to be replaced by fine connectedness, or which dimension/potential-theory conventions are intended. The cited motivation about polygonal path connectedness cannot silently remove a necessary connectedness qualification. These choices are required before writing a unique fine-potential-theoretic measure identity; no standard harmonic measure on an ordinary open domain is substituted for Keldysh measure.
\subsection{Short English statement}
Does the finite-rectifiable-path accessibility theorem for harmonic measure extend to fine potential theory and Keldysh measure? The source invokes polygonal path connectedness as motivation, but omits the precise fine-domain, boundary, path and measure conventions. The supplied triage reports openness in Hayman's 2018 update and no definitive resolution located; no independent current-status review is asserted.
\subsection{Sources}
\begin{itemize}
\item UnsolvedMath permanent ID 2307050, source code catalogue entry, ``Research Problems in Function Theory --- Problem 7.50.'' Registry curation and metadata: Ulam AI and UnsolvedMath Contributors, CC BY 4.0; updated 2026-08-06.
\item Hayman, \emph{Research Problems in Function Theory} (2018 source update), Problem 7.50, attributed to T. Lyons and B. Øksendal: \url{https://arxiv.org/abs/1809.07200}. Supplied source-TeX extraction accessed 2026-07-29. The metadata retains the update for status triage and marks status/release and rights review NEEDS\_REVIEW.
\item Brelot and Choquet are named for the ordinary-domain harmonic-measure theorem; no complete reference is supplied for it, the fine path-connectedness motivation, or Keldysh measure. No missing citation is inferred.
\end{itemize}
\end{document}
